%! AP Statistics — Unit 2, Lesson 2.7 — Warm-Up
\documentclass[10pt]{article}
\usepackage{apstats-article}
\usepackage{apstats-boxes}

\begin{document}

\pageheader{Unit 2, Lesson 2.7}{Warm-Up}
\namedateperiod

\vspace{0.3in}

% ── PART 1 ────────────────────────────────────────────────────────────────────
{\large\bfseries\color{navy} Part 1 \quad Using a Regression Equation (Review from Lesson 2.6)}

\vspace{0.12in}
\noindent A park ranger collects data on elevation above sea level (in thousands of feet, $x$)
and average daily high temperature ($^\circ$F, $y$) at several weather stations. Technology
produces the following regression equation:
\[
  \widehat{\text{temp}} = 75 - 4(\text{elevation})
\]
The data were collected at stations with elevations ranging from \textbf{1 to 8 thousand feet}.

\vspace{0.12in}

\begin{enumerate}[leftmargin=1.6em, itemsep=0.22in]

  \item What is the slope? \blank{3cm} \quad Interpret the slope in context.\\[8pt]
        \writeline\\[2pt]\writeline\\[2pt]\writeline

  \item What is the $y$-intercept? \blank{3cm} \quad Is it meaningful in context? \blank{2cm}

        Explain.\\[8pt]
        \writeline\\[2pt]\writeline

  \item Predict the average daily high temperature at a station with an elevation of
        \textbf{4 thousand feet}. Show your work.\\[8pt]
        \writeline\\[2pt]\writeline

\end{enumerate}

\vspace{0.22in}

% ── PART 2 ────────────────────────────────────────────────────────────────────
{\large\bfseries\color{navy} Part 2 \quad Setting Up for Residuals}

\vspace{0.12in}

\begin{tcolorbox}[colback=greenbg, colframe=greenacc, boxrule=0.8pt, arc=2mm,
    left=4mm, right=4mm, top=3mm, bottom=3mm]
    \small
    \textbf{Think about it:} The regression line predicts a temperature at elevation $= 4$.
    But the actual recorded temperature at that station was \textbf{57°F}.
\end{tcolorbox}

\vspace{0.14in}

\begin{enumerate}[leftmargin=1.6em, itemsep=0.22in, start=4]

  \item How far off was the prediction? Compute: actual $-$ predicted.\\[8pt]
        \writeline\\[2pt]\writeline

  \item Is the actual temperature higher or lower than predicted? What does that tell you
        about where this data point falls relative to the regression line (above or below)?\\[8pt]
        \writeline\\[2pt]\writeline

\end{enumerate}

\end{document}
